% pruned-general-case-algorithm.tex — \input'd from proofs.tex
% Conventions: thesis/CLAUDE.md
%
% Sources: Jörn's dictation + HK2017 paper. Notation per correspondence.tex.
%
% Structure: Facet adjacency graph, adjacency constraint, adjacency pruning

% Jörn: text approved (0839595) — from \begin{definition} to \end{definition}
\begin{definition}[Facet adjacency graph]\label{def:adjacency-graph}
% QC: two facets are adjacent iff they share at least one point.
%   In 4D, adjacent facets (3-faces) can share a 2-face (ridge),
%   1-face (edge), or 0-face (vertex). This uses the weakest condition.
% Downstream: Rust: precompute adjacency matrix once per polytope.
%   Test: for the 4-simplex (5 facets), all pairs are adjacent (K_5).
The \emph{facet adjacency graph} $G_K$ of a polytope $K$ has
vertex set $\{1, \ldots, F\}$ (one vertex per facet) and an edge
$\{i, j\}$ whenever $F_i \cap F_j \neq \emptyset$.
\end{definition}

% Jörn: text approved (0839595) — from \begin{lemma} to \end{proof}
\begin{lemma}[Adjacency constraint]\label{lem:adjacency-constraint}
% QC: follows from the generalized Reeb orbit condition + continuity.
%   Key step: velocity R_i requires the orbit to lie on F_i
%   (otherwise R_i is not available in the convex hull).
%   At the breakpoint, the orbit lies on both adjacent facets.
% Downstream: prune (S, sigma) where consecutive facets in sigma
%   are not adjacent in G_K. Significant speedup for polytopes
%   with sparse adjacency graphs.
Let $\gamma$ be a simple Reeb orbit on a polytope $K$ with
representation $(\sigma, \tau)$
(Remark~\ref{rem:simple-orbit-data}).
If $\tau_k > 0$ and $\tau_{k+1} > 0$ (indices cyclic), then
$\sigma(k)$ and $\sigma(k+1)$ are adjacent in~$G_K$:
the breakpoint $\gamma(t^*)$ lies in
$F_{\sigma(k)} \cap F_{\sigma(k+1)}$.
\end{lemma}

\begin{proof}
During interval~$k$, the orbit has velocity
$\dot{\gamma}(t) = R_{\sigma(k)}$ and therefore lies on~$F_{\sigma(k)}$
(Definition~\ref{def:reeb-orbit-polytope}: for $R_i$ to appear
in the velocity, $\gamma(t) \in F_i$ is required).
Since $F_{\sigma(k)}$ is closed and $\gamma$ is continuous,
$\gamma(t^*) \in F_{\sigma(k)}$ at the breakpoint~$t^*$.
By the same argument applied to interval~$k{+}1$,
$\gamma(t^*) \in F_{\sigma(k+1)}$.
\end{proof}

% Jörn: text approved (0839595) — from \begin{corollary} to \end{corollary}
\begin{corollary}[Adjacency pruning]\label{cor:adjacency-pruning}
% QC: in the algorithm, all facets in S have beta_i > 0 (by the filter),
%   so all are active. The adjacency constraint applies to all consecutive
%   pairs in sigma (cyclically). Equivalently: sigma must trace a
%   Hamiltonian cycle in G_K[S].
% Downstream: before solving the linear system, check adjacency of all
%   consecutive pairs in sigma. Skip (S, sigma) if any pair is non-adjacent.
%   Further: skip S entirely if G_K[S] has no Hamiltonian cycle.
In the algorithm (Section~\ref{sec:main-result}), discard any
$(S, \sigma)$ for which some consecutive pair
$\sigma(k), \sigma(k{+}1)$ (cyclically) satisfies
$F_{\sigma(k)} \cap F_{\sigma(k+1)} = \emptyset$.
Equivalently, restrict to cyclic orderings~$\sigma$ of~$S$ that
form a Hamiltonian cycle in the induced subgraph~$G_K[S]$.
\end{corollary}

